\documentclass{zhvt-classic}

% 使用本地源流明体（GenRyuMin2 TW）按文件路径加载，不依赖系统字体名
% 版框（内框／界栏／书口侧栏／鱼尾）改用与版心印章同一印泥棕色：
%   印章图片 yshs.png 实色 = RGB(135,68,52) = HTML 874434
\usepackage{xcolor}
\usepackage{tikz}
% sealbrown 双注册：xcolor 供 TikZ/pgf，L3 供 cls 版框
\definecolor{sealbrown}{HTML}{874434}
\ExplSyntaxOn
\color_set:nnn { sealbrown } { HTML } { 874434 }
\ExplSyntaxOff

\zhvtset{
  grid color       = sealbrown,
  main font        = GenRyuMin2TW-SB,
  main font path   = C:/Users/Administrator/texmf/fonts/opentype/public/GenRyuMin2TW/,
jiazhu font      = GenRyuMin2TW-M,
  jiazhu font path = C:/Users/Administrator/texmf/fonts/opentype/public/GenRyuMin2TW/,
  % 页码框容纳 4 字：正常页「九十七」+前置空 1 格=4 行恰好；
  % 100+ 页「一百零一」共 5 行超出，末字「一/二/三」被框裁掉（2026-10-04 实测），故放宽到 5
  shukou page number chars = 5,
  % 书口书名框高从默认 6 字减到 4 字，把省出的 2 字高拨给 chapter 框，
  % 使 6 字书名+卷名放进 9 字高的 chapter 框后不再压线（2026-10-05）
  shukou title chars = 4,
}

\title{史记}[平津侯主父列传]
\author{汉·司马迁}
\maker{闲叟漫记·重製}


% ===== 索隐述赞样式（2026-11-08 四次定稿 · 愚大人亲手改定参数）=====
% ⚠ 竖排前提：自定义字体家族必须带 RawFeature={vertical:+vert}；
%   simhei.ttf 加 +vhal 会整体加载失败静默回退，只写 +vert；simfang 同理。
\newCJKfontfamily{\shuzanfont}[Path=C:/Windows/Fonts/,RawFeature={vertical:+vert}]{simhei.ttf}
\newCJKfontfamily{\shuzanbody}[Path=C:/Windows/Fonts/,RawFeature={vertical:+vert}]{simfang.ttf}
% 题头：152x30pt 圆角 6pt 棕底白字；baseline (0,2) 使整框落入界行格子不压版框线；
% 白字节点 x 不补偿、y +8.2pt（像素级墨迹实测 dx/dy≈0，L/R 留白约 20pt）。
\newcommand{\shuzanhead}[1]{%
  \begin{tikzpicture}[baseline={(0,2)}]
    \fill[sealbrown, rounded corners=6pt]
      (-76pt,-15pt) rectangle (76pt,15pt);
    \node[text=white, inner sep=0pt, anchor=center] at (0pt,8.2pt)
      {\shuzanfont\fontsize{23}{23}\selectfont #1};
  \end{tikzpicture}%
}
% ★ 修正书口卷名位置：chapter 框起始位置改回 cls 默认值，避免全局影响。
%   只在每个卷首页通过 \juan 命令里的空命令微调该页位置（安全，不会跨页污染）。
\ExplSyntaxOn
\AtBeginDocument{
  \dim_set:Nn \l__zhvt_shukou_chapter_title_start_dim
    { \l__zhvt_shukou_title_height_dim + \l__zhvt_font_size_dim * 2 }
  \dim_set:Nn \l__zhvt_shukou_page_number_start_dim
    { \l__zhvt_shukou_chapter_title_start_dim + \l__zhvt_shukou_chapter_title_height_dim }
}
\ExplSyntaxOff

\begin{document}

% ★★ 版心（书口）排布覆写：把书名放到鱼尾之下（2026-10-04）
%
%   cls 默认版心自上而下：书名（title 静态框，在鱼尾之上）→ 鱼尾 → 篇名（chapter 动态框）→ 页码
%   本模板改为：          鱼尾（位置不动）→ 书名 → 篇名 → 页码分割线 → 页码
%
%   手法：把 cls 放在 title 静态框里的书名，改放进 chapter 动态框（鱼尾之下）；
%         title 静态框留空。书名取 \l__zhvt_title_str，即 \title{...} 的第一参数，
%         所以换书时只改 \title，不必在此重复写书名。
%
%   ⚠ 三条必须遵守的写法（2026-10-04 实测）：
%     ① 汉字必须用 \hbox{...} 包住。flowfram 会把框内容装进一个宽 14.4mm（= 一个 line_sep
%        = 书口栏宽）的纵向 minipage 并按该宽度折行，而整页靠 PDF /Rotate 90 显示、
%        框内文字本身不变向 —— 不包 \hbox 时「围炉夜话」4 字会被逐字折成 4 行，
%        在版面上就成了横排 4 字并横向溢出版框。
%     ② 竖排下**一个 \hbox 就是一整列**。书名与篇名要同处一列，必须放进同一个 \hbox；
%        分成两个 \hbox 时，后一个会被当成新的一列、甩到隔壁栏去
%        （实测「序」跑到书口栏外侧 14.4mm 的正文列上）。中间的 \quad 在框内是横向空白，
%        显示出来正好是书名与篇名之间的纵向间隔。
%     ③ ★ 内容盒要写成 **与框等宽的定宽盒** `\hbox to \l__zhvt_shukou_chapter_title_height_dim{…\hfil}`。
%        实测：chapter 框是把内容盒沿纵向**居中**摆放的，内容宽度一变，书名就跟着漂：
%        自然宽度的 \hbox（书名 38.4mm）落在 x=157.2mm，加篇名后（57.6mm）落在 x=152.4mm，
%        目录页（无篇名）与正文页的书名就差半格。做成等宽盒后，无论有没有篇名，
%        内容固定从框顶起排，实测两页均落在 x=142.8mm，不再漂。
\ExplSyntaxOn
% 把书口 chapter 框（书名+篇名）起始位置下移 1 字高，避免 4 字卷名顶到鱼尾
% ⚠ cls 的 \__zhvt_calculate_dimensions: 在 \AtBeginDocument 里执行，会重算该 dim；
%   必须把覆写挂到 begin document 之后，否则会被二次覆盖（2026-10-05 实测）。
% (旧方案已移到导言区，直接修改 chapter 框起始位置)
% 卷标题命令 #1=书口篇名(鱼尾) #2=正文卷名（含卷序，卷名内空格用 ~）
%
% ★★ 目录页两项修订（2026-10-05，愚大人指出）★★
%   ① 页码起点统一对齐：
%      原状 —— 目录项写成 \par\hspace{4em}卷名\hspace{4em}页码，页码起点 = 卷名实宽 + 4em。
%      卷名实宽随字数变：「卷一~梁惠王上」6 字 +~、「卷七~离娄上」5 字 +~，
%      差 1 个字位 → 卷七~卷十 4 条的页码整体比其余 10 条高 27.21pt（一个字），看着上下不齐。
%      改法 —— 卷名补全角空格（宽 = 1 字 = 27.215pt）补足到 **7 字符位**再写进目录，
%      14 条页码起点遂全部落在同一条竖线上（起点对齐；末字随页码字数自然参差）。
%      ⚠ 只对目录里的副本 \l_tmpc_tl 补，正文卷名仍用 \l_tmpa_tl 原样 —— 正文像素不受影响。
%   ② 卷名整体上移半字位：前导 \hspace 由 4em 改 3.5em。
%      实测原状目录块在版框内偏下（天头留白 139.6pt / 地脚留白 110.7pt），
%      上移 0.5em = 13.6pt 后上下基本均分，视觉居中。
\NewDocumentCommand { \juan } { mm } {
	\clearpage
	\tl_set:Nn \l__zhvt_chapter_title_tl { #1 }
	\tl_set:Nn \l_tmpa_tl { #2 }
	\tl_set_eq:NN \l_tmpc_tl \l_tmpa_tl
	\int_step_inline:nn { 7 - \str_count:n { #2 } }
	  { \tl_put_right:Nn \l_tmpc_tl { 　 } }
	\tl_set:Nx \l_tmpb_tl { \zhnum { page } }
	\tl_put_left:Nn \l_tmpb_tl { \hspace { 4em } }
	\tl_put_left:Nn \l_tmpb_tl { \l_tmpc_tl }
	\tl_put_left:Nn \l_tmpb_tl { \par \hspace { 3.5em } }
	\addtocontents { toc } { \l_tmpb_tl }
	\skip_horizontal:n { \l__zhvt_font_size_dim * 4 }
	\l_tmpa_tl
	\par
	\thispagestyle { empty }
}
\cs_set_protected:Npn \__zhvt_set_shukou_content: {
	% ★ 鱼尾上方留白（title 静态框）嵌入闲叟钤印：
	%   整页 /Rotate 90，故此处把图用 \rotatebox{90} 转正为竖章；
	%   尺寸卡在版心条幅内（书内约 7mm 宽 × 33.8mm 高，留余量不压鱼尾）。
	\setstaticcontents * { title } {
		\vbox_to_ht:nn { \l__zhvt_baseline_skip_dim } {
			\vskip 14.2pt
			\hbox_to_wd:nn { \l__zhvt_shukou_title_height_dim } {
				\hfil
				\rotatebox [ origin = c ] { 90 } { \includegraphics [ width = 7mm , height = 33.8mm ] { yshs.png } }
				\hfil
			}
		}
	}
% ★ 版心（书口）统一改法（2026-11-16 愚大人定）：18pt＋框内居中＋书口文字同版框色
%   ① 原 27.2pt 时 7 字长篇名（如「廉颇蔺相如列传」）末字溢出 chapter 框、压住页码；
%      改 \fontsize{18pt}{18pt}，前后 \hfil 令字串在框内居中（不随篇名字数漂移）。
%   ② \vbox{\vskip 8.2pt} 定量下移，抵消缩字号后栏宽方向 8.2pt 上偏（像素实测 y 中心 543.3 = 原版）。
%      ⚠ \raisebox / \vspace* 在 flowfram 框内无效；换字号须同步换补偿：16pt→10.1、19pt→7.4。
%   ③ chapter 与 page~number 均置 \color{sealbrown}（HTML 874434），与版框／界行同色；正文仍黑。
%   ⚠ 页码字位、钤印、鱼尾、界行、正文一律未动；五部合订本同此口径。
	\setdynamiccontents * { chapter }
	{ \color{sealbrown} \vbox { \vskip 8.2pt \hbox to \l__zhvt_shukou_chapter_title_height_dim
		{ \hfil \fontsize{18pt}{18pt}\selectfont \l__zhvt_title_str \quad \l__zhvt_chapter_title_tl \hfil } } }
	\setdynamiccontents * { page~number } { \color{sealbrown} \quad \zhnum { page } }
}
\__zhvt_set_shukou_content:
\ExplSyntaxOff

\maketitle

% 单篇，无多卷目录，不渲染目录页

\juan{平津侯主父列传}{卷一百十二~平津侯主父列传第五十二}
\hspace*{2em}丞相公孙弘者。齐菑川国薛县人也。字季。少时为薛狱吏。有罪。免。家贫。牧豕海上。年四十馀。乃学春秋杂说。养后母孝谨。\\
\hspace*{2em}建元元年。天子初即位。招贤良文学之士。是时弘年六十。徵以贤良为博士。使匈奴。还报。不合上意。上怒。以为不能。弘乃病免归。\\
\hspace*{2em}元光五年。有诏徵文学。菑川国复推上公孙弘。弘让谢国人曰。﹁臣已尝西应命。以不能罢归。原更推选。﹂国人固推弘。弘至太常。太常令所徵儒士各对策。百馀人。弘第居下。策奏。天子擢弘对为第一。召入见。状貌甚丽。拜为博士。是时通西南夷道。置郡。巴蜀民苦之。诏使弘视之。还奏事。盛毁西南夷无所用。上不听。\\
\hspace*{2em}弘为人恢奇多闻。常称以为人主病不广大。人臣病不俭节。弘为布被。食不重肉。后母死。服丧三年。每朝会议。开陈其端。令人主自择。不肯面折庭争。於是天子察其行敦厚。辩论有馀。习文法吏事。而又缘饰以儒术。上大说之。二岁中。至左内史。弘奏事。有不可。不庭辩之。尝与主爵都尉汲黯请间。汲黯先发之。弘推其后。天子常说。所言皆听。以此日益亲贵。尝与公卿约议。至上前。皆倍其约以顺上旨。汲黯庭诘弘曰。﹁齐人多诈而无情实。始与臣等建此议。今皆倍之。不忠。﹂上问弘。弘谢曰。﹁夫知臣者以臣为忠。不知臣者以臣为不忠。﹂上然弘言。左右幸臣每毁弘。上益厚遇之。\\
\hspace*{2em}元朔三年。张欧免。以弘为御史大夫。是时通西南夷。东置沧海。北筑朔方之郡。弘数谏。以为罢敝中国以奉无用之地。原罢之。於是天子乃使硃买臣等难弘置朔方之便。发十策。弘不得一。弘乃谢曰。﹁山东鄙人。不知其便若是。原罢西南夷。沧海而专奉朔方。﹂上乃许之。\\
\hspace*{2em}汲黯曰。﹁弘位在三公。奉禄甚多。然为布被。此诈也。﹂上问弘。弘谢曰。﹁有之。夫九卿与臣善者无过黯。然今日庭诘弘。诚中弘之病。夫以三公为布被。诚饰诈欲以钓名。且臣闻管仲相齐。有三归。侈拟於君。桓公以霸。亦上僭於君。晏婴相景公。食不重肉。妾不衣丝。齐国亦治。此下比於民。今臣弘位为御史大夫。而为布被。自九卿以下至於小吏。无差。诚如汲黯言。且无汲黯忠。陛下安得闻此言。﹂天子以为谦让。愈益厚之。卒以弘为丞相。封平津侯。\\
\hspace*{2em}弘为人意忌。外宽内深。诸尝与弘有卻者。虽详与善。阴报其祸。杀主父偃。徙董仲舒於胶西。皆弘之力也。食一肉脱粟之饭。故人所善宾客。仰衣食。弘奉禄皆以给之。家无所馀。士亦以此贤之。\\
\hspace*{2em}淮南。衡山谋反。治党与方急。弘病甚。自以为无功而封。位至丞相。宜佐明主填抚国家。使人由臣子之道。今诸侯有畔逆之计。此皆宰相奉职不称。恐窃病死。无以塞责。乃上书曰。﹁臣闻天下之通道五。所以行之者三。曰君臣。父子。兄弟。夫妇。长幼之序。此五者天下之通道也。智。仁。勇。此三者天下之通德。所以行之者也。故曰﹃力行近乎仁。好问近乎智。知耻近乎勇﹄。知此三者。则知所以自治。知所以自治。然后知所以治人。天下未有不能自治而能治人者也。此百世不易之道也。今陛下躬行大孝。鉴三王。建周道。兼文武。厉贤予禄。量能授官。今臣弘罢驽之质。无汗马之劳。陛下过意擢臣弘卒伍之中。封为列侯。致位三公。臣弘行能不足以称。素有负薪之病。恐先狗马填沟壑。终无以报德塞责。原归侯印。乞骸骨。避贤者路。﹂天子报曰。﹁古者赏有功。褎有德。守成尚文。遭遇右武。未有易此者也。朕宿昔庶几获承尊位。惧不能宁。惟所与共为治者。君宜知之。盖君子善善恶恶。君若谨行。常在朕躬。君不幸罹霜露之病。何恙不已。乃上书归侯。乞骸骨。是章朕之不德也。今事少间。君其省思虑。一精神。辅以医药。﹂因赐告牛酒杂帛。居数月。病有瘳。视事。\\
\hspace*{2em}元狩二年。弘病。竟以丞相终。子度嗣为平津侯。度为山阳太守十馀岁。坐法失侯。\\
\hspace*{2em}主父偃者。齐临菑人也。学长短纵横之术。晚乃学易。春秋。百家言。游齐诸生间。莫能厚遇也。齐诸儒生相与排摈。不容於齐。家贫。假贷无所得。乃北游燕。赵。中山。皆莫能厚遇。为客甚困。孝武元光元年中。以为诸侯莫足游者。乃西入关见卫将军。卫将军数言上。上不召。资用乏。留久。诸公宾客多厌之。乃上书阙下。朝奏。暮召入见。所言九事。其八事为律令。一事谏伐匈奴。其辞曰。\\
\hspace*{2em}臣闻明主不恶切谏以博观。忠臣不敢避重诛以直谏。是故事无遗策而功流万世。今臣不敢隐忠避死以效愚计。原陛下幸赦而少察之。\\
\hspace*{2em}司马法曰。﹁国虽大。好战必亡。天下虽平。忘战必危。﹂天下既平。天子大凯。春蒐秋狝。诸侯春振旅。秋治兵。所以不忘战也。且夫怒者逆德也。兵者凶器也。争者末节也。古之人君一怒必伏尸流血。故圣王重行之。夫务战胜穷武事者。未有不悔者也。昔秦皇帝任战胜之威。蚕食天下。并吞战国。海内为一。功齐三代。务胜不休。欲攻匈奴。李斯谏曰。﹁不可。夫匈奴无城郭之居。委积之守。迁徙鸟举。难得而制也。轻兵深入。粮食必绝。踵粮以行。重不及事。得其地不足以为利也。遇其民不可役而守也。胜必杀之。非民父母也。靡弊中国。快心匈奴。非长策也。﹂秦皇帝不听。遂使蒙恬将兵攻胡。辟地千里。以河为境。地固泽卤。不生五穀。然后发天下丁男以守北河。暴兵露师十有馀年。死者不可胜数。终不能逾河而北。是岂人众不足。兵革不备哉。其势不可也。又使天下蜚刍輓粟。起於黄。腄。琅邪负海之郡。转输北河。率三十锺而致一石。男子疾耕不足於粮饟。女子纺绩不足於帷幕。百姓靡敝。孤寡老弱不能相养。道路死者相望。盖天下始畔秦也。\\
\hspace*{2em}及至高皇帝定天下。略地於边。闻匈奴聚於代谷之外而欲击之。御史成进谏曰。﹁不可。夫匈奴之性。兽聚而鸟散。从之如搏影。今以陛下盛德攻匈奴。臣窃危之。﹂高帝不听。遂北至於代谷。果有平城之围。高皇帝盖悔之甚。乃使刘敬往结和亲之约。然后天下忘干戈之事。故兵法曰﹁兴师十万。日费千金﹂。夫秦常积众暴兵数十万人。虽有覆军杀将系虏单于之功。亦適足以结怨深雠。不足以偿天下之费。夫上虚府库。下敝百姓。甘心於外国。非完事也。夫匈奴难得而制。非一世也。行盗侵驱。所以为业也。天性固然。上及虞夏殷周。固弗程督。禽兽畜之。不属为人。夫上不观虞夏殷周之统。而下近世之失。此臣之所大忧。百姓之所疾苦也。且夫兵久则变生。事苦则虑易。乃使边境之民弊靡愁苦而有离心。将吏相疑而外市。故尉佗。章邯得以成其私也。夫秦政之所以不行者。权分乎二子。此得失之效也。故周书曰﹁安危在出令。存亡在所用﹂。原陛下详察之。少加意而熟虑焉。\\
\hspace*{2em}是时赵人徐乐。齐人严安俱上书言世务。各一事。徐乐曰。臣闻天下之患在於土崩。不在於瓦解。古今一也。何谓土崩。秦之末世是也。陈涉无千乘之尊。尺土之地。身非王公大人名族之后。无乡曲之誉。非有孔。墨。曾子之贤。陶硃。猗顿之富也。然起穷巷。奋棘矜。偏袒大呼而天下从风。此其故何也。由民困而主不恤。下怨而上不知。俗已乱而政不脩。此三者陈涉之所以为资也。是之谓土崩。故曰天下之患在於土崩。何谓瓦解。吴。楚。齐。赵之兵是也。七国谋为大逆。号皆称万乘之君。带甲数十万。威足以严其境内。财足以劝其士民。然不能西攘尺寸之地而身为禽於中原者。此其故何也。非权轻於匹夫而兵弱於陈涉也。当是之时。先帝之德泽未衰而安土乐俗之民众。故诸侯无境外之助。此之谓瓦解。故曰天下之患不在瓦解。由是观之。天下诚有土崩之势。虽布衣穷处之士或首恶而危海内。陈涉是也。况三晋之君或存乎。天下虽未有大治也。诚能无土崩之势。虽有彊国劲兵不得旋踵而身为禽矣。吴。楚。齐。赵是也。况群臣百姓能为乱乎哉。此二体者。安危之明要也。贤主所留意而深察也。\\
\hspace*{2em}间者关东五穀不登。年岁未复。民多穷困。重之以边境之事。推数循理而观之。则民且有不安其处者矣。不安故易动。易动者。土崩之势也。故贤主独观万化之原。明於安危之机。脩之庙堂之上。而销未形之患。其要。期使天下无土崩之势而已矣。故虽有彊国劲兵。陛下逐走兽。射蜚鸟。弘游燕之囿。淫纵恣之观。极驰骋之乐。自若也。金石丝竹之声不绝於耳。帷帐之私俳优侏儒之笑不乏於前。而天下无宿忧。名何必汤武。俗何必成康。虽然。臣窃以为陛下天然之圣。宽仁之资。而诚以天下为务。则汤武之名不难侔。而成康之俗可复兴也。此二体者立。然后处尊安之实。扬名广誉於当世。亲天下而服四夷。馀恩遗德为数世隆。南面负扆摄袂而揖王公。此陛下之所服也。臣闻图王不成。其敝足以安。安则陛下何求而不得。何为而不成。何征而不服乎哉。严安上书曰。\\
\hspace*{2em}臣闻周有天下。其治三百馀岁。成康其隆也。刑错四十馀年而不用。及其衰也。亦三百馀岁。故五伯更起。五伯者。常佐天子兴利除害。诛暴禁邪。匡正海内。以尊天子。五伯既没。贤圣莫续。天子孤弱。号令不行。诸侯恣行。彊陵弱。众暴寡。田常篡齐。六卿分晋。并为战国。此民之始苦也。於是彊国务攻。弱国备守。合从连横。驰车击毂。介胄生虮虱。民无所告愬。\\
\hspace*{2em}及至秦王。蚕食天下。并吞战国。称号曰皇帝。主海内之政。坏诸侯之城。销其兵。铸以为锺虡。示不复用。元元黎民得免於战国。逢明天子。人人自以为更生。乡使秦缓其刑罚。薄赋敛。省繇役。贵仁义。贱权利。上笃厚。下智巧。变风易俗。化於海内。则世世必安矣。秦不行是风而其故俗。为智巧权利者进。笃厚忠信者退。法严政峻。谄谀者众。日闻其美。意广心轶。欲肆威海外。乃使蒙恬将兵以北攻胡。辟地进境。戍於北河。蜚刍輓粟以随其后。又使尉屠睢将楼船之士南攻百越。使监禄凿渠运粮。深入越。越人遁逃。旷日持久。粮食绝乏。越人击之。秦兵大败。秦乃使尉佗将卒以戍越。当是时。秦祸北构於胡。南挂於越。宿兵无用之地。进而不得退。行十馀年。丁男被甲。丁女转输。苦不聊生。自经於道树。死者相望。及秦皇帝崩。天下大叛。陈胜。吴广举陈。武臣。张耳举赵。项梁举吴。田儋举齐。景驹举郢。周市举魏。韩广举燕。穷山通谷豪士并起。不可胜载也。然皆非公侯之后。非长官之吏也。无尺寸之势。起闾巷。杖棘矜。应时而皆动。不谋而俱起。不约而同会。壤长地进。至于霸王。时教使然也。秦贵为天子。富有天下。灭世绝祀者。穷兵之祸也。故周失之弱。秦失之彊。不变之患也。\\
\hspace*{2em}今欲招南夷。朝夜郎。降羌僰。略濊州。建城邑。深入匈奴。燔其茏城。议者美之。此人臣之利也。非天下之长策也。今中国无狗吠之惊。而外累於远方之备。靡敝国家。非所以子民也。行无穷之欲。甘心快意。结怨於匈奴。非所以安边也。祸结而不解。兵休而复起。近者愁苦。远者惊骇。非所以持久也。今天下锻甲砥剑。桥箭累弦。转输运粮。未见休时。此天下之所共忧也。夫兵久而变起。事烦而虑生。今外郡之地或几千里。列城数十。形束壤制。旁胁诸侯。非公室之利也。上观齐晋之所以亡者。公室卑削。六卿大盛也。下观秦之所以灭者。严法刻深。欲大无穷也。今郡守之权。非特六卿之重也。地几千里。非特闾巷之资也。甲兵器械。非特棘矜之用也。以遭万世之变。则不可称讳也。\\
\hspace*{2em}书奏天子。天子召见三人。谓曰。﹁公等皆安在。何相见之晚也。﹂於是上乃拜主父偃。徐乐。严安为郎中。数见。上疏言事。诏拜偃为谒者。迁为中大夫。一岁中四迁偃。\\
\hspace*{2em}偃说上曰。﹁古者诸侯不过百里。彊弱之形易制。今诸侯或连城数十。地方千里。缓则骄奢易为淫乱。急则阻其彊而合从以逆京师。今以法割削之。则逆节萌起。前日晁错是也。今诸侯子弟或十数。而適嗣代立。馀虽骨肉。无尺寸地封。则仁孝之道不宣。原陛下令诸侯得推恩分子弟。以地侯之。彼人人喜得所原。上以德施。实分其国。不削而稍弱矣。﹂於是上从其计。又说上曰。﹁茂陵初立。天下豪桀并兼之家。乱众之民。皆可徙茂陵。内实京师。外销奸猾。此所谓不诛而害除。﹂上又从其计。\\
\hspace*{2em}尊立卫皇后。及发燕王定国阴事。盖偃有功焉。大臣皆畏其口。赂遗累千金。人或说偃曰。﹁太横矣。﹂主父曰。﹁臣结发游学四十馀年。身不得遂。亲不以为子。昆弟不收。宾客弃我。我阸日久矣。且丈夫生不五鼎食。死即五鼎烹耳。吾日暮途远。故倒行暴施之。﹂\\
\hspace*{2em}偃盛言朔方地肥饶。外阻河。蒙恬城之以逐匈奴。内省转输戍漕。广中国。灭胡之本也。上览其说。下公卿议。皆言不便。公孙弘曰。﹁秦时常发三十万众筑北河。终不可就。已而弃之。﹂主父偃盛言其便。上竟用主父计。立朔方郡。\\
\hspace*{2em}元朔二年。主父言齐王内淫佚行僻。上拜主父为齐相。至齐。遍召昆弟宾客。散五百金予之。数之曰。﹁始吾贫时。昆弟不我衣食。宾客不我内门。今吾相齐。诸君迎我或千里。吾与诸君绝矣。毋复入偃之门。﹂乃使人以王与姊奸事动王。王以为终不得脱罪。恐效燕王论死。乃自杀。有司以闻。\\
\hspace*{2em}主父始为布衣时。尝游燕。赵。及其贵。发燕事。赵王恐其为国患。欲上书言其阴事。为偃居中。不敢发。及为齐相。出关。即使人上书。告言主父偃受诸侯金。以故诸侯子弟多以得封者。及齐王自杀。上闻大怒。以为主父劫其王令自杀。乃徵下吏治。主父服受诸侯金。实不劫王令自杀。上欲勿诛。是时公孙弘为御史大夫。乃言曰。﹁齐王自杀无后。国除为郡。入汉。主父偃本首恶。陛下不诛主父偃。无以谢天下。﹂乃遂族主父偃。\\
\hspace*{2em}主父方贵幸时。宾客以千数。及其族死。无一人收者。唯独洨孔车收葬之。天子后闻之。以为孔车长者也。\\
\hspace*{2em}太史公曰。公孙弘行义虽脩。然亦遇时。汉兴八十馀年矣。上方乡文学。招俊乂。以广儒墨。弘为举首。主父偃当路。诸公皆誉之。及名败身诛。士争言其恶。悲夫。\\
\hspace*{2em}太皇太后诏大司徒大司空。﹁盖闻治国之道。富民为始。富民之要。在於节俭。孝经曰﹃安上治民。莫善於礼﹄。﹃礼。与奢也宁俭﹄。昔者管仲相齐桓。霸诸侯。有九合一匡之功。而仲尼谓之不知礼。以其奢泰侈拟於君故也。夏禹卑宫室。恶衣服。后圣不循。由此言之。治之盛也。德优矣。莫高於俭。俭化俗民。则尊卑之序得。而骨肉之恩亲。争讼之原息。斯乃家给人足。刑错之本也欤。可不务哉。夫三公者。百寮之率。万民之表也。未有树直表而得曲影者也。孔子不云乎。﹃子率而正。孰敢不正﹄。﹃举善而教不能则劝﹄。维汉兴以来。股肱宰臣身行俭约。轻财重义。较然著明。未有若故丞相平津侯公孙弘者也。位在丞相而为布被。脱粟之饭。不过一肉。故人所善宾客皆分奉禄以给之。无有所馀。诚内自克约而外从制。汲黯诘之。乃闻于朝。此可谓减於制度而可施行者也。德优则行。否则止。与内奢泰而外为诡服以钓虚誉者殊科。以病乞骸骨。孝武皇帝即制曰﹃赏有功。襃有德。善善恶恶。君宜知之。其省思虑。存精神。辅以医药﹄。赐告治病。牛酒杂帛。居数月。有瘳。视事。至元狩二年。竟以善终于相位。夫知臣莫若君。此其效也。弘子度嗣爵。后为山阳太守。坐法失侯。夫表德章义。所以率俗厉化。圣王之制。不易之道也。其赐弘后子孙之次当为后者爵关内侯。食邑三百户。徵诣公车。上名尚书。朕亲临拜焉。﹂\\
\hspace*{2em}班固称曰。公孙弘。卜式。兒宽皆以鸿渐之翼困於燕雀。远迹羊豕之间。非遇其时。焉能致此位乎。是时汉兴六十馀载。海内乂安。府库充实。而四夷未宾。制度多阙。上方欲用文武。求之如弗及。始以蒲轮迎枚生。见主父而叹息。群臣慕乡。异人并出。卜式试於刍牧。弘羊擢於贾竖。卫青奋於奴仆。日磾出於降虏。斯亦曩时版筑饭牛之朋矣。汉之得人。於兹为盛。儒雅则公孙弘。董仲舒。兒宽。笃行则石建。石庆。质直则汲黯。卜式。推贤则韩安国。郑当时。定令则赵禹。张汤。文章则司马迁。相如。滑稽则东方朔。枚皋。应对则严助。硃买臣。历数则唐都。落下闳。协律则李延年。运筹则桑弘羊。奉使则张骞。苏武。将帅则卫青。霍去病。受遗则霍光。金日磾。其馀不可胜纪。是以兴造功业。制度遗文。后世莫及。孝宣承统。纂脩洪业。亦讲论六蓺。招选茂异。而萧望之。梁丘贺。夏侯胜。韦玄成。严彭祖。尹更始以儒术进。刘向。王襃以文章显。将相则张安世。赵充国。魏相。邴吉。于定国。杜延年。治民则黄霸。王成。龚遂。郑弘。邵信臣。韩延寿。尹翁归。赵广汉之属。皆有功迹见述於后。累其名臣。亦其次也。\\

% ===== 索隐述赞（司马贞《史记索隐》）· 统一样式 =====
% 题头「索隐述赞曰」五字并入反白框；正文每四句一列，定宽盒 \hfil 居列中；
% \fontsize{24}{41} 列距对齐全书界行（121+41k）；\clearpage 使述赞独占一页。
\clearpage
\shuzanhead{索隐述赞曰}\linebreak
{\shuzanbody\fontsize{24}{41}\selectfont
\hbox to \hsize{\hfil 平津巨儒。晚年始遇。外示宽俭。内怀嫉妒。\hfil}\linebreak
\hbox to \hsize{\hfil 宠备荣爵。身受肺腑。主父推恩。观时设度。\hfil}\linebreak
\hbox to \hsize{\hfil 生食五鼎。死非时蠹。\hfil}%
\par}
\end{document}
